\documentclass[11pt,a4paper]{article}

\usepackage[T1]{fontenc}
\usepackage[utf8]{inputenc}
\usepackage[english]{babel}
\usepackage{lmodern}
\usepackage{microtype}
\usepackage[a4paper,margin=20mm,headheight=15pt]{geometry}
\usepackage{graphicx}
\usepackage{booktabs}
\usepackage{tabularx}
\usepackage{longtable}
\usepackage{array}
\usepackage{xcolor}
\usepackage{amsmath}
\usepackage{amssymb}
\usepackage{caption}
\usepackage{subcaption}
\usepackage{float}
\usepackage{enumitem}
\usepackage{pdflscape}
\usepackage{fancyhdr}
\usepackage[hidelinks]{hyperref}
\usepackage{xurl}

\graphicspath{{figures/}}
\definecolor{dsblue}{HTML}{2563EB}
\definecolor{flexorange}{HTML}{EA580C}
\definecolor{flexgreen}{HTML}{16A34A}
\definecolor{slate}{HTML}{334155}
\definecolor{lightgray}{HTML}{F1F5F9}

\newcommand{\dstwr}{\textcolor{dsblue}{DS-TWR}}
\newcommand{\flextdoa}{\textcolor{flexorange}{FlexTDOA}}
\newcommand{\code}[1]{\texttt{\detokenize{#1}}}
\newcolumntype{Y}{>{\raggedright\arraybackslash}X}
\newcolumntype{R}[1]{>{\raggedleft\arraybackslash}p{#1}}

\pagestyle{fancy}
\fancyhf{}
\lhead{\small Experimental comparison: DS-TWR vs. FlexTDOA}
\rhead{\small July 26, 2026}
\cfoot{\thepage}

\hypersetup{
  pdftitle={Experimental comparison of native DS-TWR and FlexTDOA},
  pdfauthor={UWB ESP-IDF technical report},
  pdfsubject={Static LOS test, four anchors, centered tag},
  pdfkeywords={UWB, DS-TWR, FlexTDOA, DW3000, ESP32-S3}
}

\begin{document}
\pagenumbering{gobble}

\begin{titlepage}
  \centering
  \vspace*{17mm}
  {\Huge\bfseries Experimental comparison\\[3mm]
  Native \dstwr{} vs. \flextdoa\par}
  \vspace{7mm}
  {\Large Field data, surveyed $3\times3$ m geometry, static centered tag\par}
  \vspace{18mm}

  \begin{tabular}{rl}
    Test date: & July 26, 2026, 20:04--20:18 EEST\\
    Platform: & 5 ESP32-S3 + DW3000 modules, Raspberry Pi 5\\
    Firmware: & commit \texttt{a97cdce60d19}, \texttt{-dirty} state\\
    Branch: & \texttt{experiment/ds-twr-native-performance}\\
    Configuration: & M1 tag; M2, M3, M4, M5 anchors\\
    Environment: & LOS, static position\\
  \end{tabular}

  \vfill
  \colorbox{lightgray}{%
    \begin{minipage}{0.88\textwidth}
      \vspace{3mm}
      \textbf{Key result.}
      At this central test point and under the tested conditions, DS-TWR
      delivered the best unfiltered accuracy: a 2D RMSE of
      \textbf{3.25 cm} and P95 of \textbf{4.80 cm}, at
      \textbf{15.15 positions/s}. FlexTDOA delivered much higher radio
      throughput with a passive tag: \textbf{804.6 observations/s}; the
      embedded solver published \textbf{25.20 positions/s}, with a
      \textbf{3.59 cm} median, \textbf{7.56 cm} P95, and
      \textbf{4.44 cm} RMSE after the ground-truth-independent structural
      validity checks described in Section~4.3.
      \vspace{3mm}
    \end{minipage}}
  \vfill
  {\large Reproducible technical report\par}
  \vspace{5mm}
\end{titlepage}

\pagenumbering{roman}
\tableofcontents
\clearpage
\pagenumbering{arabic}

\section{Executive summary}

The test addressed two distinct questions:

\begin{enumerate}[leftmargin=7mm]
  \item how accurate the measurements and resulting position are without
  temporal filtering, in the same physical geometry;
  \item how quickly and robustly each implementation delivers useful data in
  the current configuration.
\end{enumerate}

Six alternating 120 s blocks were collected:
\texttt{Flex1--DS1--Flex2--DS2--Flex3--DS3}. Each protocol therefore has
360 s of observation time, after 15 s of stabilization following each mode
change. No dashboard polling errors occurred. The final data set contains
22,256 DS-TWR ranges, 289,716 captured FlexTDOA observations, 5,455 complete
DS positions, 50,879 offline FlexTDOA solutions from complete valid slots, and
9,075 positions published by the ESP32 solver. Four FlexTDOA observations and
one embedded position failed structural validity checks and remain available
in the audit data, but are excluded from all comparative metrics.

\begin{table}[H]
  \centering
  \caption{Aggregate position results after structural validity checks. No
  position was selected or rejected using its ground-truth error.}
  \footnotesize
  \begin{tabular}{lrrrrrr}
    \toprule
    Series & $N$ & Hz & Median & P95 & RMSE & Maximum\\
    \midrule
    DS-TWR, surveyed geometry
      & 5,455 & 15.15 & 3.04 cm & 4.80 cm & 3.25 cm & 6.79 cm\\
    FlexTDOA, offline complete slot
      & 50,879 & 141.30 & 7.64 cm & 17.17 cm & 9.94 cm & 84.53 cm\\
    FlexTDOA, ESP32 solver
      & 9,074 & 25.20 & 3.59 cm & 7.56 cm & 4.44 cm & 25.44 cm\\
    \bottomrule
  \end{tabular}
\end{table}

\paragraph{Preliminary interpretation.}
At this single static point, DS-TWR produced lower position errors. FlexTDOA
provided a passive tag and higher update rate, at the cost of greater
dispersion. The result does not establish the general superiority of either
protocol: the central LOS point is geometrically favorable, and the test does
not cover edges, motion, antenna orientation, or NLOS conditions.

\section{Technical background}

\subsection{Native DS-TWR}

The tested implementation uses exactly the three required UWB frames:
\texttt{POLL}, \texttt{RESP}, and \texttt{FINAL}. The tag initiates the
exchange with each anchor; \texttt{FINAL} carries the tag timestamps, and the
anchor combines them with its own hardware timestamps to calculate the
propagation time. This is the flow described in the official Qorvo DW3xxx
examples~\cite{qorvo-api}.

For asymmetric DS-TWR, propagation time is estimated using the
relationship~\cite{aps013}:

\begin{equation}
  \widehat{T}_{\mathrm{prop}} =
  \frac{R_aR_b-D_aD_b}{R_a+D_a+R_b+D_b},
\end{equation}

where $R_a,R_b$ are the round-trip intervals and $D_a,D_b$ are the response
delays measured by the two clocks. The distance is
$\widehat d=c\widehat T_{\mathrm{prop}}$.

In the firmware, the timestamps include calibrated antenna delays, and the
estimate also uses instantaneous CIA/carrier-integrator correction. These are
physical measurement corrections, not temporal filters. Each result is a
single \texttt{POLL--RESP--FINAL} measurement; the dashboard applies no
median, moving average, Kalman filter, or smoothing before trilateration.

\subsection{FlexTDOA}

FlexTDOA is distributed DL-TDOA: within a slot, one anchor transmits
\texttt{REQ}, and $K$ anchors respond in sequence. The tag only listens. The
initiator and responder order rotate, eliminating a permanently designated
reference anchor while keeping the tag radio-passive. The architecture, CI-CR
scheduling, and AlgMin come from Pătru et al.~\cite{flextdoa-paper}.

For an observation with initiator $i$ and responder $j$, the firmware forms:

\begin{equation}
 z_{ij} =
 c\left[(t^{\mathrm{tag}}_j-t^{\mathrm{tag}}_i)
 -\widehat T_{\mathrm{reply},j}\right]
 -\widehat d_{ij},
\end{equation}

and ideally:

\begin{equation}
 z_{ij} \simeq
 \lVert\mathbf p-\mathbf a_j\rVert -
 \lVert\mathbf p-\mathbf a_i\rVert .
\end{equation}

The reply term is corrected using instantaneous CFO, and
$\widehat d_{ij}$ is the anchor-to-anchor distance carried in the protocol
stream. The AlgMin position minimizes:

\begin{equation}
  \min_{\mathbf p}
  \sum_{(i,j)}
  \left[
  \lVert\mathbf p-\mathbf a_j\rVert -
  \lVert\mathbf p-\mathbf a_i\rVert-z_{ij}
  \right]^2 .
\end{equation}

The firmware applies no median, moving average, or Kalman filtering to the
observations. The embedded solver uses the most recent still-valid directed
observations (at most 12 for four anchors) in a damped least-squares solution.
Reusing the latest coherent set is not a temporal average, but it makes the
embedded output different from the offline ``one slot, three observations''
solution reported separately here.

\begin{figure}[H]
  \centering
  \includegraphics[width=\textwidth]{09_protocol_timing.pdf}
  \caption{Timing schedule active during the test. P-R-F =
  \texttt{POLL--RESP--FINAL}.}
\end{figure}

\begin{table}[H]
  \centering
  \caption{Structural differences relevant to the system.}
  \small
  \begin{tabularx}{\textwidth}{p{31mm}YY}
    \toprule
    Characteristic & Native DS-TWR & FlexTDOA\\
    \midrule
    Tag role & Active; transmits two of the three frames for every range
      & Radio-passive; receive only\\
    Observable & Absolute tag-to-anchor distance & Range difference between
      two anchors\\
    Frames & 3 per range: POLL, RESP, FINAL & 1 REQ + up to 3 RESP per slot\\
    Scaling with tags & UWB traffic grows with the number of tags & Additional
      tags require no additional UWB slots\\
    Solver in test & Linear trilateration on the Raspberry Pi & 2D AlgMin on
      ESP32 plus offline reproduction\\
    Temporal filtering & No & No\\
    Active profile & 64 ms/frame with 4 anchors & 12 ms/frame with 4 slots\\
    \bottomrule
  \end{tabularx}
\end{table}

\section{Test configuration}

\subsection{Geometry and reference}

The anchors were surveyed as a square with a 3.000 m side and a stated maximum
positioning deviation of $\pm2$ mm. The convention used in the common analysis
is:

\[
 A2=(0,0),\quad A3=(0,3),\quad A4=(3,0),\quad A5=(3,3)\ {\rm m},
\]
\[
 \mathbf p_{\rm tag}=(1.5,\,1.5)\ {\rm m}.
\]

The true tag-to-anchor distance is the same for all anchors:

\[
 d_{\text{true}}=\sqrt{1.5^2+1.5^2}
 =2.121320344\ {\rm m}.
\]

At the center, all true TDOA differences are zero. This property allows each
$z_{ij}$ to be evaluated directly without introducing solver error into the
measurement analysis.

\begin{figure}[H]
  \centering
  \includegraphics[width=0.86\textwidth]{01_geometry.pdf}
  \caption{Surveyed geometry used for the common comparison and NVS geometry
  used by the operational solver.}
\end{figure}

The fixed geometry stored in NVS at the time of the test was:

\begin{center}
\begin{tabular}{lrrr}
  \toprule
  Anchor & NVS coordinate [m] & Surveyed coordinate [m] & Difference [mm]\\
  \midrule
  A2 & $(0.000,\ 0.000)$ & $(0.000,\ 0.000)$ & $(0,\ 0)$\\
  A3 & $(0.000,\ 3.045)$ & $(0.000,\ 3.000)$ & $(0,\ +45)$\\
  A4 & $(3.030,\ -0.058)$ & $(3.000,\ 0.000)$ & $(+30,\ -58)$\\
  A5 & $(3.040,\ 2.880)$ & $(3.000,\ 3.000)$ & $(+40,\ -120)$\\
  \bottomrule
\end{tabular}
\end{center}

Consequently, the embedded FlexTDOA positions were evaluated in their
operational coordinate frame relative to the center of the NVS geometry,
$(1.5175,\ 1.46675)$ m. For the direct protocol comparison, both offline
solutions use the surveyed square and the $(1.5,\ 1.5)$ m reference. Keeping
these evaluations separate prevents anchor self-localization error from being
attributed to the positioning protocol.

\subsection{Radio and timing parameters}

Both protocols used the same DW3000 hardware, channel 5, 6.8 Mb/s, 64 MHz PRF,
and preamble length 128. The DW3000 family is designed for TWR and TDoA, and
the manufacturer specifies ranging accuracy below 10 cm in its product
specification~\cite{qorvo-product,dwm3000-datasheet}; the result in this report,
however, applies only to the specific platform and environment tested.

\begin{table}[H]
  \centering
  \caption{Active runtime profiles read from all five modules.}
  \small
  \begin{tabularx}{\textwidth}{lR{25mm}Y}
    \toprule
    Parameter & Value & Interpretation\\
    \midrule
    DS slot & 15 ms & budget for one native exchange\\
    DS gap after 4 anchors & 4 ms & complete frame $4\times15+4=64$ ms\\
    DS RX slice / timeout & 100 / 8 ms & anchor listening / expected-frame timeout\\
    DS RESP / FINAL delay & 2 / 2 ms & hardware delayed TX\\
    DS auto-RX delay & 500 UUS & RX opening after TX\\
    \midrule
    Flex guard & 250 $\mu$s & per-slot guard\\
    Flex REQ subslot / process & 250 / 1,000 $\mu$s & request and processing\\
    Flex RESP subslot & 250 $\mu$s & spacing between responses\\
    Flex process per responder & 250 $\mu$s & $3\times250\ \mu$s\\
    Flex slot / frame & 3 / 12 ms & 333.33 slots/s; 83.33 frames/s\\
    \bottomrule
  \end{tabularx}
\end{table}

The active antenna delays, in DW3000 register units, were: M1 16,363,
M2 16,368, M3 16,362, M4 16,371, and M5 16,371.

\section{Methodology}

\subsection{Experimental plan}

\begin{itemize}[leftmargin=6mm]
  \item six alternating 120 s Flex--DS blocks, reducing confounding between
  the protocol and thermal or temporal drift;
  \item a 15 s warm-up after every mode change;
  \item dashboard polling at 50 Hz and deduplication by event ID, sequence,
  slot, and pair;
  \item full status at 1 Hz; all five modules remained online and in the
  expected mode;
  \item zero polling errors across all six blocks;
  \item the system was left in DS-TWR mode at the end.
\end{itemize}

\subsection{Three position evaluations}

\paragraph{Offline DS-TWR, surveyed geometry.}
A position is calculated only when one range for each of A2--A5 exists within
the same frame. The same unweighted linear trilateration as in the dashboard
was used. No ranges were removed based on error magnitude.

\paragraph{Offline FlexTDOA, surveyed geometry.}
A slot is solved only when the tag received all three responders. The 2D
AlgMin implementation reproduces the dashboard solver convention and
iterations, starting deterministically at the center of the square. Each
solution uses only that slot; there is no temporal seed, averaging, or
filtering.

\paragraph{Embedded FlexTDOA, operational geometry.}
The positions published directly by the ESP32-S3 are analyzed. They use up to
12 recent directed differences and the NVS geometry. This is the series
relevant to what the user sees in the dashboard.

\subsection{Structural validity rules and metrics}

The comparison uses validity checks that do not depend on the known tag
position. A corrected FlexTDOA observation is accepted only if it satisfies
the universal range-difference bound
\[
 |z_{ij}|\leq\|\mathbf a_i-\mathbf a_j\|+0.01\ {\rm m}.
\]
Four of 289,716 captured observations failed this check. An embedded position
is accepted only when the solver's own \texttt{rms\_m} does not exceed the
anchor-array diagonal, $\sqrt{3^2+3^2}=4.243$ m. One of 9,075 published
positions failed this deliberately permissive self-consistency check. DS-TWR
had no corresponding structural failure. All excluded records remain in the
derived CSV files with a validity flag and in the immutable raw captures.

No ground-truth error threshold, percentile trimming, or manual event ID is
used for selection. The same accepted data population is used consistently
for aggregate metrics, per-block metrics, and plots.

We report axis bias, median radial error, P95, P99, maximum, 2D RMSE, and
precision relative to the sample centroid (CEP50/CEP95). For measurements, we
report bias, standard deviation, MAE, absolute P95, and robust deviation
$1.4826\times{\rm MAD}$.

Repeatability is evaluated across the three blocks rather than treating
hundreds of thousands of temporally correlated samples as independent
observations. Where a 95\% uncertainty is stated, it is the Student interval
of the mean of the three blocks ($t_{0.975,2}=4.303$). With only three blocks,
the interval should be interpreted descriptively.

\section{Results}

\subsection{Throughput, availability, and position rate}

\begin{table}[H]
  \centering
  \caption{Per-block results. For Flex, ``complete'' means three distinct
  responses in a slot; for DS, four ranges in a position frame. The Flex
  position is the embedded output. Flex counts and rates shown here are after
  structural validation; 0, 2, and 2 observations and 0, 1, and 0 positions
  were excluded from Flex 1, 2, and 3, respectively.}
  \scriptsize
  \begin{tabular}{lrrrrrr}
    \toprule
    Block & Measurements & Measurements/s & Theoretical delivery & Complete &
      Positions/s & Median / P95\\
    \midrule
    Flex 1 & 95,932 & 799.07 & 79.91\% & 40.82\% & 25.66 & 3.72 / 7.56 cm\\
    DS 1   & 7,426  & 61.88  & 99.01\% & 97.70\% & 15.21 & 3.05 / 4.84 cm\\
    Flex 2 & 96,808 & 806.69 & 80.67\% & 43.01\% & 24.94 & 3.56 / 7.57 cm\\
    DS 2   & 7,400  & 61.66  & 98.66\% & 96.41\% & 15.01 & 3.05 / 4.83 cm\\
    Flex 3 & 96,972 & 807.95 & 80.80\% & 43.45\% & 25.00 & 3.49 / 7.52 cm\\
    DS 3   & 7,430  & 61.91  & 99.06\% & 97.91\% & 15.24 & 3.01 / 4.73 cm\\
    \bottomrule
  \end{tabular}
\end{table}

The block-level means were:

\begin{itemize}[leftmargin=6mm]
  \item DS-TWR: $61.82\pm0.34$ ranges/s and
  $15.15\pm0.31$ complete frames/s (95\%);
  \item FlexTDOA: $804.57\pm11.93$ valid observations/s,
  $141.30\pm11.68$ solvable complete slots/s, and
  $25.20\pm1.00$ valid embedded positions/s (95\%);
  \item the Flex observation rate was 13.0 times the DS range rate;
  the complete Flex slot rate was 9.3 times the complete DS frame rate;
  the actual embedded output was 1.66 times faster than the DS position rate.
\end{itemize}

DS delivery was 98.91\% of the profile maximum, and 97.34\% of reconstructed
candidate frames contained all four ranges. With Flex, the tag observed
80.46\% of the theoretical maximum of 1,000 differences/s; 42.43\% of slots
contained all three responders. A rolling solver can nevertheless publish
positions using recent pairs, so the percentage of complete slots is not
identical to embedded-position availability.

\begin{figure}[H]
  \centering
  \includegraphics[width=0.94\textwidth]{07_throughput_completeness.pdf}
  \caption{Observed throughput and completeness of position units. The
  throughput axes compare different observables: DS absolute ranges and Flex
  range differences.}
\end{figure}

\subsection{DS-TWR measurement quality}

\begin{table}[H]
  \centering
  \caption{DS-TWR range error relative to 2.121320344 m.}
  \small
  \begin{tabular}{lrrrrrr}
    \toprule
    Anchor & $N$ & Bias & $\sigma$ & $\sigma_{\rm MAD}$ & MAE & Absolute P95\\
    \midrule
    A2 & 5,561 & +2.45 cm & 1.77 cm & 1.78 cm & 2.56 cm & 5.37 cm\\
    A3 & 5,555 & +8.02 cm & 0.92 cm & 0.89 cm & 8.02 cm & 9.77 cm\\
    A4 & 5,571 & +3.03 cm & 1.53 cm & 1.63 cm & 3.04 cm & 5.47 cm\\
    A5 & 5,569 & $-4.21$ cm & 1.56 cm & 1.48 cm & 4.22 cm & 6.53 cm\\
    \bottomrule
  \end{tabular}
\end{table}

DS-TWR has low dispersion on each link, but its bias is module-dependent:
from $-4.21$ cm at A5 to $+8.02$ cm at A3. This indicates a clear opportunity
for per-module/antenna recalibration. The opposing biases project partly
favorably onto the central position; the position result should not be
interpreted as evidence that every range is calibrated to within 3 cm.

\begin{figure}[H]
  \centering
  \includegraphics[width=0.78\textwidth]{02_ds_range_error.pdf}
  \caption{Unfiltered distribution of DS-TWR errors. Outliers are not drawn in
  the boxplot, but they are included in every metric. The maximum absolute
  error was 11.47 cm.}
\end{figure}

\subsection{FlexTDOA observation quality}

At the center, the true value for every directed pair is zero. The CFO plus
anchor-to-anchor propagation correction is essential:

\begin{table}[H]
  \centering
  \caption{The 289,712 structurally valid FlexTDOA observations, aggregated.
  The nearly zero corrected global bias includes cancellation among directed
  biases and should not be interpreted on its own.}
  \small
  \begin{tabular}{lrrrrrr}
    \toprule
    Series & Bias & $\sigma_{\rm MAD}$ & MAE & P95 abs. & P99 abs. & Max. abs.\\
    \midrule
    Corrected & 0.02 cm & 12.31 cm & 9.94 cm & 24.10 cm & 36.40 cm &
      211.4 cm\\
    Uncorrected & 71.78 cm & 35.43 cm & 71.92 cm & 122.30 cm & 131.40 cm &
      161.2 cm\\
    \bottomrule
  \end{tabular}
\end{table}

Across pairs, corrected bias ranges from $-13.61$ to $+13.18$ cm, while robust
dispersion ranges from 7.71 to 9.34 cm. The correction drastically reduces the
normal offset but does not eliminate the distribution tails.

Exactly four corrected observations violated the physical inequality
$|z_{ij}|\leq\|\mathbf a_i-\mathbf a_j\|+1$ cm: two in Flex 2 and two in
Flex 3, all with A3 as initiator, in two isolated slots. The values were
approximately 79.6--79.7 m. They are retained in the audit data but are not
used by the measurement statistics, offline solver, or comparative plots.
The exclusion rate is 4 out of 289,716, or 0.00138\%.

\begin{figure}[H]
  \centering
  \includegraphics[width=\textwidth]{03_flex_pair_heatmaps.pdf}
  \caption{Bias and robust dispersion by directed pair after structural
  validation. The uncorrected raw field primarily contains the reply delay
  and clock offset.}
\end{figure}

\subsection{Position accuracy}

\begin{table}[H]
  \centering
  \caption{Aggregate position metrics after structural validity checks. CEP is
  calculated relative to each distribution's own centroid and describes
  precision, not absolute accuracy.}
  \scriptsize
  \begin{tabular}{lrrrrrrrr}
    \toprule
    Series & $x/y$ bias & RMSE & Median & P95 & P99 & CEP50/95 &
      $>50$ cm\\
    \midrule
    DS, surveyed geometry &
      $+2.63/-0.99$ cm & 3.25 & 3.04 & 4.80 & 5.51 &
      1.35 / 2.78 cm & 0 / 5,455\\
    Flex, offline complete slot &
      $+4.53/+0.60$ cm & 9.94 & 7.64 & 17.17 & 25.50 &
      6.44 / 15.42 cm & 53 / 50,879\\
    Flex, operational ESP32 &
      $+2.29/+0.39$ cm & 4.44 & 3.59 & 7.56 & 10.44 &
      2.84 / 6.64 cm & 0 / 9,074\\
    \bottomrule
  \end{tabular}
\end{table}

DS-TWR was more precise and more accurate in the raw comparison on the
surveyed geometry. Relative to DS, FlexTDOA solved from individual slots had
3.06 times the RMSE and 3.58 times the P95, but delivered 9.3 times as many
solvable complete slots per second.

The embedded FlexTDOA solver is significantly more stable than a single-slot
solution because it uses multiple recent pairs: P95 is 7.56 cm and RMSE is
4.44 cm. P95 is 1.58 times the DS value, while RMSE is 1.37 times the DS
value. This improvement is not the result of an explicit temporal filter; it
comes from overdetermination and reuse of the current observation set.

\begin{figure}[H]
  \centering
  \includegraphics[width=\textwidth]{04_position_clouds.pdf}
  \caption{Position clouds. The cross marks the reference for the respective
  geometric frame. All plotted positions passed structural validation.}
\end{figure}

\begin{figure}[H]
  \centering
  \includegraphics[width=0.85\textwidth]{05_position_error_cdf.pdf}
  \caption{Cumulative distribution of radial error using all samples. The axis
  limit is robust; the maximum outside the axis is annotated in the plot.}
\end{figure}

\subsection{Records excluded by structural validation}

The range-difference bound excluded four FlexTDOA observations, and the
solver-self-consistency bound excluded one embedded position. The latter had
an internal RMS of 21.854 m, more than five times the permissive 4.243 m
limit; the next-highest internal RMS in the data set was 0.617 m. Therefore,
the exclusion is determined by a large, isolated break in solver consistency,
not by its position error relative to the known center.

The raw captures and the unfiltered embedded-position CSV retain the excluded
records and their reasons. This preserves auditability while preventing an
invalid protocol state from dominating the static accuracy comparison.

\subsection{Temporal repeatability}

DS-TWR had per-block RMSE values of 3.20--3.28 cm and P95 values of
4.73--4.84 cm. Mean per-block RMSE is $3.248\pm0.095$ cm (95\%). Offline Flex
had RMSE of 9.90--10.01 cm and P95 of 16.96--17.34 cm, with a mean RMSE of
$9.940\pm0.154$ cm. Embedded Flex had RMSE of 4.39--4.50 cm, with a
mean of $4.443\pm0.127$ cm, and P95 of 7.52--7.57 cm.

\begin{figure}[H]
  \centering
  \includegraphics[width=\textwidth]{06_position_error_time.pdf}
  \caption{Median error in 1 s intervals. Individual events are not visible in
  a per-second median and are analyzed separately.}
\end{figure}

\begin{figure}[H]
  \centering
  \includegraphics[width=0.88\textwidth]{08_block_repeatability.pdf}
  \caption{Per-block RMSE and P95 for the operational output: offline DS and
  embedded Flex, after structural validation.}
\end{figure}

\section{Interpretation and protocol selection}

\subsection{When DS-TWR is the appropriate choice}

\begin{itemize}[leftmargin=6mm]
  \item one tag or few tags, where the additional traffic is acceptable;
  \item accuracy and short distribution tails have priority;
  \item it is acceptable for the tag to transmit and consume radio energy;
  \item a rate of approximately 15 positions/s meets the application's
  dynamics.
\end{itemize}

In the current test, DS has the best combination of RMSE, P95, between-block
stability, and frame completeness. Per-anchor bias nevertheless shows that
recalibration could improve accuracy further or alter it at other positions.

\subsection{When FlexTDOA is the appropriate choice}

\begin{itemize}[leftmargin=6mm]
  \item a fully UWB-passive tag, with battery life and scalability as
  priorities;
  \item multiple tags that must be located without consuming additional
  slots;
  \item a high observation rate and faster updates;
  \item the application tolerates a P95 of approximately 7.6 cm in the current
  configuration and validates impossible outputs.
\end{itemize}

FlexTDOA is not simply ``faster DS''; its observable, topology, and scaling
cost are different. With four anchors, the 12 ms profile produces
83.33 frames/s and 333.33 slots/s, independent of the number of passive tags.
For a multi-tag installation, the airtime comparison becomes progressively
more favorable to FlexTDOA.

\subsection{Immediate recommendations}

\begin{enumerate}[leftmargin=7mm]
  \item Retain unfiltered native DS-TWR with the 64 ms profile as the accuracy
  baseline.
  \item Recalibrate the antenna delays, especially for the A3 and A5 links,
  and then repeat the test exactly.
  \item Write the surveyed geometry to NVS before operational comparisons;
  the current geometry differs by up to 12 cm on the $y$ axis.
  \item For FlexTDOA, add the physical gate
  $|z_{ij}|\leq d_{ij}+\varepsilon$ and a quality gate on RMS/bounding box.
  These are validity checks, not temporal filtering.
  \item Repeat the test on a grid of at least $5\times5$ points, including the
  edges and outside of the square, followed by dynamic and NLOS testing. The
  center alone does not characterize GDOP over the full area.
  \item For motion, report end-to-end latency using hardware timestamps, not
  only the rate of events received by the dashboard.
\end{enumerate}

\section{Limitations}

\begin{itemize}[leftmargin=6mm]
  \item A single, static, central point cannot describe spatial error.
  \item The test was LOS; NLOS and severe multipath are not represented.
  \item The geometry is stated to $\pm2$ mm, but no external optical system or
  total station was included in the digital capture.
  \item Temperature and RSSI were recorded but not varied in a controlled way.
  \item The DS position is calculated on the Raspberry Pi, while the
  operational Flex position is calculated on the ESP32; the report therefore
  also includes a common offline Flex solution.
  \item Throughput observed through telemetry/dashboard is an end-to-end
  implementation limit, not only the protocol's PHY capacity.
  \item The data set comes from \texttt{-dirty} firmware; the manifest and
  commit provide identification, but exact reproduction requires preservation
  of the local changes that were active.
\end{itemize}

\section{Preliminary findings}

For the surveyed $3\times3$ m configuration with a centered tag under LOS:

\begin{itemize}[leftmargin=6mm]
  \item \textbf{DS-TWR produced lower static position error at this point}:
  RMSE 3.25 cm, median 3.04 cm, P95 4.80 cm, with no errors above 6.79 cm
  among complete positions;
  \item \textbf{FlexTDOA provided higher throughput and inherent multi-tag
  scalability}: 804.6
  observations/s, a radio-passive tag, and 25.2 embedded positions/s;
  \item \textbf{Embedded FlexTDOA remained close to DS-TWR in its accepted
  population}, with a 3.59 cm median, 7.56 cm P95, and 4.44 cm RMSE;
  \item \textbf{the validity policy had negligible data-volume impact}: four
  of 289,716 Flex observations and one of 9,075 embedded positions were
  excluded without consulting ground-truth position error;
  \item \textbf{the product determines the choice}: DS for accuracy with few
  tags; FlexTDOA for passive and numerous tags at a high update rate.
\end{itemize}

These findings are repeatable for the measured static central point because of
the alternating blocks, large data volume, and between-block stability. They
are not a general protocol ranking. A spatial LOS/NLOS grid, dynamic trials,
and antenna-orientation trials are required before making a system-level
choice.

\clearpage
\appendix

\begin{landscape}
\section{FlexTDOA metrics by directed pair}

\scriptsize
\begin{longtable}{lrrrrrrr}
  \caption{Structurally valid corrected FlexTDOA observations; the true value
  is 0 m for the centered tag. The final column records source observations
  excluded by the physical bound.
  $\sigma_{\rm MAD}=1.4826\times{\rm MAD}$.}\\
  \toprule
  Pair & $N$ & Bias [cm] & $\sigma_{\rm MAD}$ [cm] & P95 abs. [cm] &
  P99 abs. [cm] & Max. abs. [cm] & Excluded by bound\\
  \midrule
  \endfirsthead
  \toprule
  Pair & $N$ & Bias [cm] & $\sigma_{\rm MAD}$ [cm] & P95 abs. [cm] &
  P99 abs. [cm] & Max. abs. [cm] & Excluded by bound\\
  \midrule
  \endhead
  A2$\to$A3 & 25,583 & +4.92 & 7.86 & 20.00 & 30.60 & 101.60 & 0\\
  A2$\to$A4 & 25,192 & +0.89 & 9.04 & 20.00 & 31.21 & 108.60 & 0\\
  A2$\to$A5 & 24,354 & $-8.80$ & 9.19 & 26.40 & 52.75 & 211.40 & 0\\
  A3$\to$A2 & 23,516 & $-5.71$ & 7.71 & 21.20 & 36.20 & 111.60 & 1\\
  A3$\to$A4 & 22,517 & $-7.64$ & 8.45 & 23.10 & 34.20 & 82.50 & 1\\
  A3$\to$A5 & 20,732 & $-13.61$ & 7.86 & 27.90 & 37.80 & 87.10 & 2\\
  A4$\to$A2 & 24,892 & $-0.69$ & 8.45 & 18.90 & 30.80 & 93.60 & 0\\
  A4$\to$A3 & 24,316 & +7.18 & 8.45 & 23.10 & 35.49 & 104.40 & 0\\
  A4$\to$A5 & 23,612 & $-5.98$ & 8.75 & 22.10 & 33.30 & 100.80 & 0\\
  A5$\to$A2 & 24,976 & +6.73 & 9.04 & 24.90 & 42.20 & 177.40 & 0\\
  A5$\to$A3 & 24,931 & +13.18 & 8.30 & 28.55 & 37.77 & 124.10 & 0\\
  A5$\to$A4 & 25,091 & +5.87 & 9.34 & 23.20 & 34.50 & 107.80 & 0\\
  \bottomrule
\end{longtable}
\end{landscape}

\section{Files and reproduction}

The artifacts are stored in:
\begin{quote}
\path{/home/pi/Documents/UWB/reports/uwb_protocol_comparison_20260726}
\end{quote}

\small
\noindent
\begin{tabularx}{\textwidth}{p{56mm}Y}
  \toprule
  File/directory & Contents\\
  \midrule
  \path{data/*.jsonl.xz} & lossless raw-capture archives, one file per block\\
  \path{data/*.metadata.json} & configuration, status, and counters for each
  block\\
  \path{raw_data_manifest.csv} & size and SHA-256 for every raw and compressed
  file\\
  \path{analysis_summary.json} & all aggregate metrics\\
  \path{measurement_metrics.csv} & per-anchor/pair metrics\\
  \path{position_metrics.csv} & aggregate position metrics\\
  \path{block_metrics.csv} & results for all six blocks\\
  \path{ds_positions.csv} & reconstructed DS positions\\
  \path{flex_slot_positions.csv} & Flex positions from complete slots\\
  \path{flex_embedded_positions.csv} & positions published by the ESP32\\
  \path{figures/} & all plots in vector PDF and PNG formats\\
  \bottomrule
\end{tabularx}

\paragraph{Reproducing the analysis}
\begin{verbatim}
cd /home/pi/Documents/UWB/reports/uwb_protocol_comparison_20260726
xz --decompress --keep data/*.jsonl.xz
cd /home/pi/Documents/UWB
python3 tools/uwb_compare_analyze.py
cd reports/uwb_protocol_comparison_20260726
latexmk -pdf -interaction=nonstopmode -halt-on-error report.tex
\end{verbatim}

The collection script is \path{tools/uwb_compare_collect.py}, and the analysis
script is \path{tools/uwb_compare_analyze.py}. The analysis does not modify the
raw files.

\begin{thebibliography}{9}

\bibitem{flextdoa-paper}
G.-C. Pătru, L. Flueratoru, I. Vasilescu, D. Niculescu, and D. Rosner,
``FlexTDOA: Robust and Scalable Time-Difference of Arrival Localization Using
Ultra-Wideband Devices,'' \emph{IEEE Access}, vol. 11, 2023,
pp. 28610--28627.
\href{https://doi.org/10.1109/ACCESS.2023.3259320}
{doi:10.1109/ACCESS.2023.3259320}.

\bibitem{qorvo-api}
Qorvo, \emph{DW3xxx Device Driver API Guide}, Sections 6.3.26--6.3.29,
double-sided two-way ranging examples, 2025.
\url{https://forum.qorvo.com/uploads/short-url/xD3TlXKvkujjdUaJWXv2b4E7GQN.pdf}.

\bibitem{aps013}
Decawave, \emph{APS013: DW1000 and Two-Way Ranging}, rev. 2.3,
Poll--Response--Final and the asymmetric DS-TWR formula.
\url{https://forum.qorvo.com/uploads/short-url/x34DrF7EW5fQP9wY3aNESqPKz8z.pdf}.

\bibitem{qorvo-product}
Qorvo, \emph{DW3110 Ultra-Wideband Transceiver}, DW3000 specifications and
documentation.
\url{https://www.qorvo.com/products/p/DW3110}.

\bibitem{dwm3000-datasheet}
Qorvo, \emph{DWM3000 IEEE 802.15.4-z UWB Transceiver Module Data Sheet},
rev. B, May 2021. Local copy:
\path{datasheets/DWM3000 Data Sheet.pdf}.

\bibitem{firmware}
Local repository \emph{UWB Module ESP-IDF}, commit
\texttt{a97cdce60d191360300f8f39e6afa77247a59f77}, branch
\texttt{experiment/ds-twr-native-performance}, working state from
July 26, 2026.

\end{thebibliography}

\end{document}
